% Generated draft: requires mathematical and source review.
\documentclass[11pt]{article}
\usepackage[margin=25mm]{geometry}
\usepackage{fontspec}
\setmainfont{Latin Modern Roman}
\usepackage{amsmath,amssymb,mathtools,unicode-math}
\setmathfont{Latin Modern Math}
\usepackage{enumitem,needspace,xurl,hyperref}
\setlength{\parindent}{0pt}
\setlength{\parskip}{5pt}
\setlength{\emergencystretch}{3em}
\newcommand{\N}{\mathbb{N}}
\newcommand{\Z}{\mathbb{Z}}
\newcommand{\Q}{\mathbb{Q}}
\newcommand{\R}{\mathbb{R}}
\newcommand{\C}{\mathbb{C}}
\newcommand{\F}{\mathbb{F}}
\newcommand{\A}{\mathbb{A}}
\newcommand{\PP}{\mathbb{P}}
\newcommand{\E}{\mathbb{E}}
\newcommand{\eps}{\varepsilon}
\newcommand{\dd}{\,\mathrm{d}}
\begin{document}
\section*{UnsolvedMath ID: 30006432}
\textbf{Source problem code:} catalogue entry
\subsection{Definitions and mathematical statement}
Let \(F/\Q\) be an arbitrary finite extension, and let \(\mathcal O_F\) denote its ring of integers. For a nonzero integral ideal \(\mathfrak a\subseteq\mathcal O_F\), put \(N\mathfrak a=\#(\mathcal O_F/\mathfrak a)\). The Dedekind zeta function is
\[
\zeta_F(s)=\sum_{0\ne\mathfrak a\subseteq\mathcal O_F}(N\mathfrak a)^{-s},\qquad \operatorname{Re}(s)>1.
\]
Use the standard Euler--Kronecker normalization: writing the Laurent expansion at \(s=1\) as
\[
\zeta_F(s)=\frac{\kappa_F}{s-1}+c_F+O(s-1),\qquad \kappa_F>0,
\]
define
\[
\gamma_F:=\frac{c_F}{\kappa_F}
=\lim_{s\to1^+}\left(\frac{\zeta_F'(s)}{\zeta_F(s)}+\frac{1}{s-1}\right)\in\R.
\]
This normalization gives
\[
\gamma_{\Q}=\gamma
=\lim_{n\to\infty}\left(\sum_{k=1}^{n}\frac1k-\log n\right).
\]
The supplied original extract stops before the definition of \(\gamma_F\); the standard normalization above makes the intended invariant explicit and requires confirmation against the complete report.

\textbf{Research question.} For an arbitrary number field \(F\), determine what arithmetic information about \(F\) is encoded by \(\gamma_F\), and characterize that information through prime ideals and their norms, zeros of \(\zeta_F\), class-field data, or other intrinsic arithmetic invariants of \(F\).

\textbf{Unresolved scope.} The supplied question does not specify what constitutes information being encoded, which invariants are to be recovered or related, or what form a satisfactory characterization must take. It therefore specifies a research direction rather than a uniquely determined conjecture with a precise conclusion. The original extract also contains an incomplete reference to Euler--Kronecker constants along \(\Z_p\)-extensions; it does not specify a prime \(p\), an extension, its layers, or a proposed assertion. No restriction to such extensions is imposed here, because the supplied cleaned statement concerns arbitrary number fields.
\subsection{Short English statement}
For each number field \(F\), the Euler--Kronecker constant \(\gamma_F\) is the ratio of the constant Laurent coefficient of its Dedekind zeta function at \(1\) to the residue there. The question asks for an arithmetic interpretation of this real number. The intended characterization remains unspecified, and a fragment mentioning \(\Z_p\)-extensions leaves a possible additional scope unresolved.
\subsection{Sources}
\begin{itemize}
\item Local UnsolvedMath registry: \texttt{data/unsolvedmath/problems.json}, permanent ID \(30006432\), source problem code \texttt{catalogue entry}, title ``Arithmetic Meaning of Euler-Kronecker Constants,'' in the supplied Oberwolfach Reports Open Problems collection. The supplied record labels the problem open and its statement \texttt{corrected\_verified}; these are registry labels, not independent mathematical verification.
\item Supplied report citation: \emph{Cohomology Theories for Automorphic Forms and Enumerative Algebra} (2026), Oberwolfach Reports, DOI \texttt{10.4171/owr/2025/38}, pp.\ 2065--2124. \url{https://doi.org/10.4171/owr/2025/38}. The record describes the report as documenting workshop open-problem discussions. It gives proposed year 2025 and does not identify a proposer or a problem-specific page.
\item Supplied literature source labeled ``Ihara primary paper'': \url{https://www.kurims.kyoto-u.ac.jp/~kenkyubu/emeritus/ihara/Publications-and-Recent-Preprints/RecentArticles/pdf-files/1A.pdf}. The original extract attributes the introduction of the invariant \(\gamma_F\) to Ihara, referring to an unidentified local reference [1].
\item Supplied literature source labeled ``Ghosh higher constants'': \url{https://doi.org/10.1142/S1793042126500521}.
\item Supplied literature source labeled ``Oberwolfach Report 38'': \url{https://ems.press/content/serial-article-files/52251}.
\item Literature assessment dated exactly \texttt{2026-08-22}: the supplied curation reports explicit formulas and links to small prime norms, zero-free regions, and bounds, while stating that no complete arithmetic interpretation is known and that the extracted question remains open. This is supplied dated curation derived from the OWR collection, not an independently verified literature assessment. The supplied background explicitly calls for independent verification of time-sensitive status and cited claims.
\end{itemize}
\end{document}
